\documentclass[10pt,a4paper]{amsart}
\usepackage[margin=19mm]{geometry}
\usepackage{amsmath,amssymb,mathtools}
\usepackage[scr=rsfs,cal=euler]{mathalfa}
\usepackage{xfrac}
\usepackage[unicode,hidelinks,bookmarks=false]{hyperref}
\newcommand{\ie}{i.e.\ }
\newcommand{\cf}{cf.\ }
\newcommand{\resp}{respectively\ }
\newcommand{\Subsection}{\subsection}
\providecommand{\nobreakdash}{\nobreak\hbox{-}}
\setlength{\emergencystretch}{2em}
\multlinegap0pt
\usepackage{bm}
\usepackage{bbm}
\usepackage{dsfont}
\usepackage{xspace}
\usepackage{cancel}
\usepackage{mathtools}
\usepackage{amsthm}
\makeatletter
\@ifundefined{them}{\newtheorem{them}{Them}}{}
\@ifundefined{cor}{\newtheorem{cor}[them]{Corollary}}{}
\@ifundefined{lem}{\newtheorem{lem}[them]{Lemma}}{}
\makeatother
\newcommand{\QQ}{\mbox{$\mathcal{Q}_\ell$}}
\newcommand{\X}{\mathcal{X}}
\newcommand{\Y}{\mathcal{Y}}
\title[Left Ehresmann monoids with a proper basis]{Left Ehresmann monoids with a proper basis}
\author{Gracinda Gomes, Victoria Gould, Yanhui Wang}
\date{Source: arXiv:2601.22923v2 (CC BY 4.0)}
\begin{document}
\maketitle

\noindent\emph{Source and licence.} Left Ehresmann monoids with a proper basis. Version arXiv:2601.22923v2 (CC BY 4.0), distributed under the Creative Commons Attribution 4.0 International licence, as recorded in the arXiv licence metadata for this exact version and shown on the abstract page https://arxiv.org/abs/2601.22923v2. Full rights notice: licenses/arxiv-2601.22923-CC-BY-4.0.txt.\par\smallskip

\noindent\textit{Editorial scope.} Counted unit: the Them whose source label is \texttt{them:main} in the article (source0.tex lines 1602--1605, source label \texttt{them:main}), delivered together with the article's own complete original proof of it (source0.tex lines 1607--1663, one \texttt{proof} environment of 57 lines). No mathematics is written, repaired or re-derived: statement and proof are the article's own text line for line. Required context retained uncounted: h+m-Hormal form (lines 1000--1008); prop (lines 1078--1087); lem:htrick (lines 1554--1559); lem:theta (lines 1575--1578); cor:nf (lines 1588--1590). Every definition, display and label of those lines is retained unchanged. Omitted from this package and not redistributed: 1--999, 1009--1077, 1088--1553, 1560--1574, 1579--1587, 1591--1601, 1606--1606, 1664--1812. Nothing inside a retained excerpt is omitted or altered: every retained body is delivered exactly as the article's lines are written, so no omission receipt of any kind accompanies this package. Citations: the retained excerpts carry no citation command at all, so no bibliography entry of the article is transcribed and no reference is redistributed. Reference adaptation: none was needed and none was made; every reference inside the delivered unit resolves inside it, so no unresolved reference or citation is produced. Printed slips kept literal: no printed word, symbol, hypothesis or number of the counted statement or of its proof was altered. Layout and class: the article's own class is \texttt{amsart} with packages including tikz, amssymb, a4, mathrsfs, cancel, pstricks, pst-node, pst-plot, verbatim, pgf, pifont, bbding, inputenc, graphicx; the wrapper is the lane's pinned \texttt{amsart} editorial wrapper at 10pt on A4, which restates the theorem-like environments the retained text needs and adds the article's own macro definitions that the retained text needs (\texttt{\textbackslash QQ}, \texttt{\textbackslash X}, \texttt{\textbackslash Y}), with extra wrapper packages (bm, bbm, dsfont, xspace, cancel, mathtools). The delivered numbering is the unit's own. Assets: no retained span contains a figure environment, a graphics-inclusion command, a PDF-page inclusion, a table or a TikZ picture; no image file is copied into this package, the package declares no asset dependency and no image file is redistributed with it. Licence: version arXiv:2601.22923v2 of this article is distributed under the Creative Commons Attribution 4.0 International licence, as recorded in the arXiv licence metadata for this exact version and shown on the abstract page https://arxiv.org/abs/2601.22923v2. A full rights notice is delivered at licenses/arxiv-2601.22923-CC-BY-4.0.txt. Characters: every retained excerpt is pure ASCII, so the delivered engine, XeLaTeX with its default Latin Modern text font, reports no missing character.
\begin{lem}\label{h+m-Hormal form}
 Let $M=\langle H\rangle_{(2)}$  be a left Ehresmann monoid  where  $H$ is a  pre-atomic $*$-subset. 
 Suppose that $h \in H$ and $m = h_1h_2\dots h_n\in M$ is in $H$-canonical form. Then $hm$ can be written as an expression in $H$-canonical form as follows:
$$ hm=
	\begin{cases}
        (hh_1)h_2\dots h_n & \mbox{if\ either}\ h_1\in E\mbox{ or}\ (h_1\notin E\ \text{and}\ h^*h_1^+= h_1^+)\\
        (hh_1^+)h_1\dots h_n & \mbox{if}\  h_1\notin E\mbox{ and } h^*h_1^+< h_1^+.\\
	\end{cases}$$
\end{lem}

\begin{lem}\label{prop}\label{prop:extra} 
Let $M=\langle H\rangle_{(2)}$ be a left Ehresmann monoid with 
pre-basis $H$. Then the following conditions hold:

\begin{enumerate}
    \item[$(i)$] if  $h_1\dots h_n$ is  in $H$-canonical form, then $(h_1\dots h_n)^+=h_1^+$.
   
    \item[$(ii)$]  if $h \in H$ and $k \in H\setminus E$, then  $h^{\ast} \geq k^{+}$ if and only if  $hk\in H$, and then $(hk)^\ast = k^\ast$.
\end{enumerate}
\end{lem}

\begin{lem}\label{lem:htrick} We have that
\[H^{\mathcal{Q}}=\{ [h](1,h^*)^\omega:h\in H\},\]
and for $h\in H,$ 
\[([h](1,h^*)^\omega)^+=h^+\kappa\mbox{ and }
([h](1,h^*)^\omega)^*=h^*\kappa.\]
\end{lem}

\begin{lem}\label{lem:theta} Define $\theta:H\rightarrow H^{\mathcal{Q}}$ by  $h\theta=[h](1,h^*)^\omega$.
Then $\theta$ is a bijection such that 
$h^*\theta=(h\theta)^*$ and $h^+\theta=(h\theta)^+$.
\end{lem}

\begin{cor}\label{cor:nf} An element $h_1\dots h_n$ of $Q$ is in $H$-canonical form if and only if 
$(h_1\theta)\dots (h_n\theta)$ is in $H^{\mathcal{Q}}$-canonical form.
\end{cor}

\begin{them}\label{them:main}   Let $Q$
be a left Ehresmann monoid with proper  basis $H$. Then $Q$ is isomorphic to  a  biunary monoid subsemigroup  $\QQ(T,\X,\Y)$ of some $\mathcal{P}_{\ell}(T,\X)$. 

\end{them}

\begin{proof} Let  $\QQ=\mathcal{Q}(T, \X , \Y)$ and $\theta:H\rightarrow H^{\mathcal{Q}}$ be the bijection defined as above. We extend the domain of $\theta$ to $Q$ by  defining 
\[(h_1\dots h_n)\theta=(h_1\theta)\dots (h_n\theta)\]
where $h_1\dots h_n$ is in $H$-canonical form. Lemma~\ref{lem:theta}  and Corollary~\ref{cor:nf} together with the fact that $H$ is a basis for $Q$ and $H^{\mathcal{Q}}$ is a basis for $\QQ$ give  us that this extended $\theta$ is also a bijection.



To show that $\theta$ is a  semigroup morphism, we proceed by induction to show  that $(h_1\dots h_n)\theta=(h_1\theta)\dots (h_n\theta)$ for any $h_i \in H$, $i = 1, \dots, n$, without the restriction that $h_1\dots h_n$ is in $H$-canonical form.

  Clearly, the statement  holds for $n =1$.  We now prove the case $n=2$, that is, we verify $(h_1h_2)\theta = (h_1\theta)( h_2\theta)$ for any $h_1,h_2\in H$.  By Lemma~\ref{h+m-Hormal form}, there are three  different possibilities when looking for the $H$-canonical form for $h_1h_2$

  (1) Suppose that $h_2\in E$. Then $[h_2]$ is the identity of $\QQ$,  and  by (H2)
we have $(h_1h_2)^* = h_1^*h_2= h_1^* h_2^*$ . Notice that  $$(1, h_1^*)^\omega(1, h_2^*)^\omega =(1, h_1^*h_2^*)^\omega$$ 
since $\kappa$ is a semilattice embedding. 
Hence
\[
\begin{aligned}
(h_1\theta)( h_2\theta) &= [h_1](1, h_1^*)^\omega[h_2](1, h_2^*)^\omega=[h_1](1, h_1^*)^\omega(1, h_2^*)^\omega=[h_1](1, h_1^*h_2^*)^\omega\\
        &=[h_1][h_2](1, h_1^*h_2^*)^\omega=[h_1h_2](1,(h_1h_2)^*)^\omega=(h_1h_2)\theta.\
\end{aligned}
\]


(2)
If $h_2\notin E$  and $h_1^* h_2^+= h_2^+$  then $h_1h_2 \in H$ and $(h_1h_2)^* = h_2^*$ by (H3). In the following have in mind that by Lemma~\ref{lem:htrick}  we  get $([h_2](1, h_2^*)^\omega)^+ = (1,h_2^+)^\omega$, and $\kappa$ is a semilattice embedding. We have   
\[
\begin{aligned}
(h_1h_2)\theta & = [h_1h_2](1, (h_1h_2)^*)^\omega
               =[h_1][h_2](1, h_2^*)^\omega=[h_1]([h_2](1, h_2^*)^\omega)^+[h_2](1, h_2^*)^\omega\\
               &=[h_1](1,h_2^+)^\omega[h_2](1, h_2^*)^\omega=[h_1](1,h_1^*)^\omega(1,h_2^+)^\omega[h_2](1, h_2^*)^\omega\\
               &=[h_1](1,h_1^*)^\omega[h_2](1, h_2^*)^\omega=(h_1\theta )(h_2\theta).
\end{aligned}
\]

(3)  If $h_2 \notin E$ and  $h_1^* h_2^+< h_2^+$,  then  Lemma~\ref{h+m-Hormal form} tells that  $(h_1h_2^+)h_2$ is in $H$-canonical form.  Using  the definition of $\theta$ and the case (1),   we have  \[(h_1h_2)\theta = ((h_1h_2^+)h_2)\theta =((h_1h_2^+)\theta )(h_2\theta)=(h_1\theta)(h_2^+\theta)(h_2\theta)\]
and then 
by Lemma~\ref{lem:theta},
\[(h_1h_2)\theta=(h_1\theta)(h_2\theta)^+(h_2\theta)=(h_1\theta)(h_2\theta).\]

Suppose for induction that  the statement holds  for $n-1$, where $n\geq 2$,  and consider a product
$h_1\dots h_n$, where $h_i\in H$ for $1\leq i\leq n$. 
Let $h_2\dots h_n = k_2\dots k_m$ where $k_2\dots k_m$  is in $H$-canonical form. Then $(h_1h_2\dots h_n)\theta = (h_1k_2\dots k_m)\theta$ and by the induction hypothesis,  $(h_2\dots h_n)\theta=(h_2\theta)\dots (h_n\theta)= (k_2\theta) \dots (k_m\theta)$. Now we look the $H$-canonical form of $h_1h_2\dots h_n=h_1k_2\dots k_m$. 
 
By Lemma~\ref{h+m-Hormal form}, this can be:

(1) $(h_1k_2)k_3\dots k_m$, in which case  
$(h_1h_2\dots h_n)\theta=  (h_1k_2)\theta (k_3\theta) \dots (k_m\theta)=(h_1\theta)( k_2\theta) (k_3\theta) \dots (k_m\theta)=(h_1\theta)(h_2\theta) \dots (h_n\theta)$ using the  $n=2$ case  and we are done too; or finally


(2) $(h_1k_2^+)k_2\dots k_m$, in which case 
$(h_1h_2\dots h_n)\theta=(h_1k_2^+)\theta (k_2\theta )\dots( k_m\theta)=(h_1\theta)(k_2^+\theta) (k_2\theta )\dots( k_m\theta)$ using the  $n=2$ case again. But $k_2^+\theta =(k_2\theta)^+$ so that $(h_1h_2\dots h_n)\theta= (h_1\theta) (k_2\theta )\dots( k_m\theta)=(h_1\theta)(h_2\theta) \dots (h_n\theta)$ as required.

It is now clear that $\theta$ is a semigroup isomorphism and as such  must preserve the identity. To show that $\theta$ respects the unary operation $+$, let $h_1h_2\dots h_n\in Q$ be in $H$-canonical form. By Lemma~\ref{prop} $(i)$, Lemma~\ref{lem:theta} and Corollary~\ref{cor:nf}, we have
\[
(h_1h_2\dots h_n)^+\theta= h_1^+\theta=(h_1\theta)^+=((h_1\theta)( h_2\theta) \dots (h_n\theta))^+= ((h_1h_2\dots h_n)\theta)^+.
\]
The argument that  $\theta$ respects the unary operation $*$ is similar.
\end{proof}

\end{document}
